\scanpage{40}
On conclut du premier groupe
\[
\varphi_*\Lambda_Y\varphi^*
=\underbrace{\varphi_*\varphi^*\Lambda_X}_{\substack{L_X\Lambda_X\\=\\\operatorname{id}_X-\sum_{0\leq i\leq\unsure{2n|n}}p_i^X}}
\left(\operatorname{id}_X-\sum_{n+1\leq j\leq2n}p_X^j\right)
=\underbrace{\operatorname{id}_X-\sum_{0\leq j\leq2n}p_X^j}_{q}.
\]
[Plus gén.
\[
\varphi_*\Lambda_Y^m\varphi^*
=\deleted{\unsure{\sum}}\left(\operatorname{id}-\sum_{0\leq i\leq n}p_X^i\right)\Lambda_X^{m-1}
\left(\operatorname{id}-\sum_{n+1\leq j\leq2n}p_X^{\unsure{i|j}}\right)
=q\Lambda_X^{m-1}q
\]
où $q=\operatorname{id}-\displaystyle\sum_{0\leq i\leq2n}p_X^i$.
\figure{\begin{tikzpicture}[x=1cm,y=1cm]\draw (0,.3)..controls(1,.5)and(2,.5)..(2.2,.15)..controls(2.3,-.1)and(2.25,-.3)..(2.6,-.3)--(7,-.4);\end{tikzpicture}}

Aussi :
\[
\varphi^*\Lambda_X\varphi_*=\operatorname{id}_Y-\underbrace{p_E}_{\text{projection sur partie évanescente }E^{n-1}}
\]
d'où, en multipliant par $\Lambda_Y$ à gauche, comme $\Lambda_Yp_E=0$
\[
\Lambda_Y=\underbrace{\Lambda_Y\varphi^*}\Lambda_X\varphi_*
=\underbrace{\varphi^*\Lambda_X(\operatorname{id}_X-\sum_{n+1\leq j\leq2n}\overbrace{p_X^j)\Lambda_X}^{\substack{0\\\Vert}}}\varphi_*
\]
donc
\[\boxed{\Lambda_Y=\varphi^*\Lambda_X^2\varphi_*}\]
Donc $C(X)$ \deleted{\unsure{équivaut}}\inserted{implique} $C(Y)$.

Soit $p_X^{i,j}$ les projections sur $L_X^jP^i(X)$ $(i+j\leq n,\ i,j\geq0)$.

Alors
\[
\varphi_*p_Y^{i,j}\inserted{\unsure{L|\Lambda}}\varphi^*=p_X^{i,j+1}
\qquad\text{pour }i+j\unsure{<|\leq}n-1,\quad i,j\geq0\quad\deleted{\unsure{algébrique|quelconque}}
\]
\margin{$L?$}
donc la condition $D(Y)$ implique que les $p_X^{i,j}$ sont algébriques pour $\unsure{\forall}\ i+j<n$, $i\geq0$ et \uline{$j\geq1$}, i.e. les seuls \unsure{rest. en susp. sont les} \uline{$p_X^i$ pour $0\leq i\leq\unsure{2n}$}. Mais
\mspage{39}
